\chapter{Station Inverse Data}
\label{chap:StatInvDat}

The Station Inverse dialog box is primarily intended to provide the user with various metrics that relate the states of two user-selected stations.  Many of these relative quantities have been defined previously: the azimuth angle in section \ref{subsection:AZIM}, the slant distance in section \ref{subsection:DIST}, the 3-D XYZ vector in section \ref{subsection:DXYZ}, the height difference in section \ref{subsection:HDIF}, and the vertical angle in section \ref{subsection:VANG}.  Note that the Station Inverse dialog box only functions after an adjustment has been run (including the calculation of any derived points), as the values within the Station Inverse dialog box require the h5 file generated by the adjustment process.  

\section{Station Inverse Dialog Data Generation}
The Station Data Dialog box shown in Figure \ref{fig:StationInverseDialog} provides a user interface through which the user can add and remove station inverse pairs. Individual information for each station is provided in the left half of the dialog box. Relative information between stations is provided in the right half of the dialog box. The center table serves as a graphical display of the station inverse pairs that will have their metrics generated in \verb+<Project Name>.inv+ and \verb+<Project Name>Inv.csv+ files. The following subsections detail the various actions the user can perform in the Station Data Dialog box and the steps that are executed to generate the output inverse pair data.

\begin{figure} [!htb]
	\centering
	\includegraphics[trim = 0mm 0mm 0mm 0mm, clip, width=1.0\textwidth]{StationInverseDialog.png} % or linewidth
	% Trim option parameter order: left bottom right top
	\caption{Station Inverse Dialog}
	\label{fig:StationInverseDialog}
\end{figure}

\subsection{Interacting with the Station Data Dialog Interface} \label{sect:SDDInter}
Upon selecting View $\ra$ Station Inverse from the main window, the Station Data Dialog box table will be loaded with all prior existing pairs. For more information regarding the parsing of the \verb+<Project Name>Inverses.cfg+ file refer to the appendix "Inverse Configuration Parsing". This interface allows the user to perform the following actions.

\begin{itemize}
	\item \emph{Add Pair to Project} - Comboboxes allow the user to user to select stations that were enabled for the last adjustment. Once a station has been selected in the To and From comboboxes they may be added to the list by clicking the Add to Project button. Note that if a station pair already exists in the list or if the To and From comboboxes have the same value, then the pair will not be added to the list.
	\item \emph{Remove Pair(s) from Project} - Left clicking a row will highlight an individual pair. Additional pairs can be highlighted individually by pressing Shift and Click on the other desired row. Additionally, once a row is selected the user can hit Ctrl + Shift + Click to highlight everything between the original row and the newly selected row. Once the desired row(s) are highlighted they may be removed by pressing Right Click and selecting Delete or by pressing the Delete key.
	\item \emph{Export Results} - Upon clicking the Export Results button all of the valid pairs in the center table will have their data generated to an output file. For more information please refer to section \ref{sect:InverseGen}
	\item \emph{Export Results As} - Behaves in a similar manner to \emph{Export Results} but with extended functionality. This button will bring up a file dialog box where the user can select the output directory to save their inv and csv files and input a new base name for the output files as well. This button should be utilized if it is not desired that the output be placed in the project directory or an output file convention differing from the default \verb+<Project Name>.inv+ and \verb+<Project Name>Inv.csv+ is desired. This allows the user to backup the current results of the inverse calculations such that they will not be overwritten (as \verb+<Project Name>.inv+ and \verb+<Project Name>Inv.csv+ will be) by subsequent modifications and adjustments to the project.
	\item \emph{Pair Selection} - Highlighting an individual pair by clicking its row, or selecting an up or down arrow will update the text boxes to display the metrics for that station inverse pair. If the pair is invalid, the boxes will be cleared.
\end{itemize}

\subsection{Inverse Pair Data Generation} \label{sect:InverseGen}
Upon clicking one of the Export buttons from the Station Data Dialog box or selecting Project $\ra$ Calculate Adjustment from the main Salsa interface the following sequence will be executed.

\begin{enumerate}
	\item Project will search for \verb+<Project Name>Inverses.cfg+ file. If not found an early exit will occur.
	\item The file will be parsed for all unique, valid pairs. A pair is considered valid if
	\begin{itemize}
		\item The From and To stations have labels that correspond to an existing POSG, POSC or derived point record.
		\item Both stations were enabled at the time Calculate Adjustment was last performed.
		\item The pair is non-trivial, i.e. the From and To stations do not have the same label. 
	\end{itemize}
	\item Output \verb+<Project Name>.inv+ and \verb+<Project Name>Inv.csv+ files will be generated containing the metrics for valid inverse pairs.
\end{enumerate}

The inv output file is a plain text file which writes each quantity from the Station Data Dialog GUI on its own line. It is meant to serve as a human readable format for all of the station inverse pairs. Conversely, the csv file delimits quantities using commas and has only one line per station pair. 

\section{Station Inverse Dialog Relative Geodetic (WGS84) Data Definitions}

As described above, all relative quantities between the two provided stations are defined in sections \ref{subsection:AZIM} through \ref{subsection:VANG}, with the exception of the ellipsoidal distance and the horizontal distance.  For these quantities previously defined, the corresponding uncertainties are also computed via the equivalent measurement first order partial derivative matrices that are employed in an adjustment.  These uncertainties are displayed within the Station Inverse dialog box next to the nominal values.  All quantities are specified in meters, except for the azimuth and vertical angles between the stations which are provided in degrees, minutes, and seconds of arc, along with the associated uncertainties in seconds of arc.  Also note that the DXYZ vector covariance is provided, which includes both the diagonal and off-diagonal covariance values. 

The ellipsoidal distance is the shortest distance along the surface of the WGS 84 ellipsoid from the ``From'' station latitude and longitude to the ``To'' station latitude and longitude.  The horizontal distance is the magnitude of the relative station vector projected into the horizontal plane, which is defined using the East-North-Up frame computed at the ``From'' station.  The ellipsoidal distance and the horizontal distance between a ``From'' and ``To'' station are shown in Figure \ref{fig:RelativeStationDistances}, along with the previously defined slant distance (section \ref{subsection:DIST}).
\begin{figure} [!htb]
	\centering
	\includegraphics[trim = 30mm 55mm 35mm 25mm, clip, width=1.0\textwidth]{DistanceDiagramV2.pdf} % or linewidth
	% Trim option parameter order: left bottom right top
	\caption{Relative Station Distances}
	\label{fig:RelativeStationDistances}
\end{figure}


\subsection{Ellipsoidal Distance}

The Ellipsoidal distance is in general defined as the shortest distance along the surface of a three dimensional ellipsoid between two points on the ellipsoid.  The ellipsoid used in SALSA is  WGS 84, which is an ellipsoid of revolution with a flattening parameter $f = 1/298.257223563$ and semi-major axis (SMA) $a_{WGS84} = 6378137.0$ meters.  The two points on the surface are often defined using their geodetic latitude and longitude values.

\subsubsection{Vincenty's Method for Ellipsoidal Distance Computation} \label{sect:Vincenty}

To compute the ellipsoidal distance between the ``From'' and ``To'' station locations, the commonly used Vincenty's Method is employed \cite{vincenty1975direct}. Vincenty's Method is an iterative numerical approach that provides consistent geodesic distance values between all points on an ellipsoid, except those points antipodal (i.e. on exactly opposite sides of the Earth) to each other (or nearly antipodal).

The semi-minor axis is needed for Vincentys method, and is computed as
\begin{equation}
b = (1 - f)a
\end{equation}

The geodesic latitude $\phi$ and longitude $\lambda$ for the two stations are assumed to be known, and the longitude difference is computed as 
\begin{equation}
\Delta\lambda = \lambda_2 - \lambda_1
\end{equation}
where $\lambda_2$ represents the longitude of the ``To'' station, and $\lambda_1$ represents the longitude of the ``From'' station.  These subscripts will be used for these stations for the remainder of this derivation.  The ``reduced latitude'' of each station (i.e. the latitude on the auxiliary sphere) is computed as
\begin{equation}
\begin{split}
U_1 &= \tan^{-1}\left[ (1-f) \tan(\phi_1) \right] \\ 
U_2 &= \tan^{-1}\left[ (1-f) \tan(\phi_2) \right]
\end{split}
\end{equation}

The quantity $\bar{\lambda}$ is initialized as the longitude difference:
\begin{equation}
\bar{\lambda} = \Delta\lambda
\end{equation}
The following equations are then iterated upon until $\bar{\lambda}$ converges (as dictated by a change in magnitude $||\Delta\bar{\lambda}||$ of less than $1 \times 10^{-12}$):
\begin{equation}
\sin \sigma = \sqrt{(\cos U_2\sin \bar{\lambda})^2 + (\cos U_1 \sin U_2 - \sin U_1 \cos U_2 \cos \bar{\lambda} )^2}
\end{equation}
\begin{equation}
\cos \sigma = \sin U_1 \sin U_2 + \cos U_1 \cos U_2 \cos \bar{\lambda}
\end{equation}
\begin{equation}
\sigma = \tan^{-1} \frac{\sin \sigma}{\cos \sigma}
\end{equation}
\begin{equation}
\sin \alpha = \frac{\cos U_1 \cos U_2 \sin \bar{\lambda}}{\sin \sigma}
\end{equation}
\begin{equation}
\cos^2 \alpha = 1 - \sin^2 \alpha
\end{equation}
\begin{equation}
\cos(2 \sigma_m) = \cos \sigma - \frac{2 \sin U_1 \sin U_2}{\cos^2 \alpha}
\end{equation}
\begin{equation}
C = \frac{f}{16} \cos^2 \alpha \left[ 4 + f \left(4 - 3 \cos^2 \alpha \right) \right]
\end{equation}
\begin{equation}
\begin{split}
\bar{\lambda}_{\text{new}} =& \Delta\lambda + (1 - C) f \sin \alpha \\ 
& * \left\lbrace \sigma + C \sin \sigma \left[ \cos(2\sigma_m) + C \cos\sigma \left(-1 + 2 \cos^2(2\sigma_m) \right) \right]  \right\rbrace 
\end{split}
\end{equation}
\begin{equation}
\Delta\bar{\lambda} = \bar{\lambda}_{\text{new}} - \bar{\lambda}_{\text{prev}}
\end{equation}

After $\bar{\lambda}$ has converged, the ellipsoidal distance $s$ is computed using equations:
\begin{equation}
u^2 = \cos^2 \alpha \frac{a^2 - b^2}{b^2}
\end{equation}
\begin{equation}
k_1 = \frac{\sqrt{1 + u^2} - 1}{\sqrt{1 + u^2} + 1} 
\end{equation}
\begin{equation}
A = \frac{1 + \tfrac{1}{4} k_1^2}{1-k_1}
\end{equation}
\begin{equation}
 B = k_1 \left(1 - \tfrac{3}{8} k_1^2 \right)
\end{equation}
\begin{equation}
\begin{split}
\Delta\sigma =& B \sin\sigma \left\lbrace \cos(2 \sigma_m) + \tfrac{1}{4} B \left[ \cos\sigma \left(-1 + 2 \cos^2(2\sigma_m) \right) \right. \right. \\
& \left. \left. - \tfrac{1}{6} B \cos(2\sigma_m) (-3 + 4 \sin^2\sigma) (-3 + 4 * \cos^2(2\sigma_m)) \right] \right\rbrace
\end{split}
\end{equation}
\begin{equation}
s = b A (\sigma - \Delta\sigma)
\end{equation}

For two points that are exactly on the equator (i.e. the geodetic latitude is zero), the value for the $\cos(2 \sigma_m)$ term is set to zero to avoid a singularity in the computation of that parameter (the term is not not used anyways when both points are exactly on the equator).  

\subsubsection{Sigma Point Transform for Uncertainty Mapping}

The unscaled sigma point transform \cite{julier1995new} is employed to map the uncertainty in both station positions into a single uncertainty value for the ellipsoidal distance. Compared to the standard linearized approach of pre- and post-multiplying the solution covariance by first order partial derivative matrices, the sigma point transform better captures the performance of nonlinear functions, while eliminating the need to compute the partial derivative matrices.

The two station positions in ECEF XYZ coordinates are stacked into a single vector as
\begin{equation}
\boldsymbol{x} = \left[ \begin{array}{c} 
\boldsymbol{r}_{\text{From}} \\ \boldsymbol{r}_{\text{To}}
\end{array} \right]
\end{equation}

The ECEF XYZ covariance of both stations is $\text{Cov}_x$, which includes cross covariance terms between the stations. To obtain $\text{Cov}_x$, the covariance values associated with these two stations are extracted from the total covariance matrix:
\begin{equation}
\text{Cov}_x = \left[ \begin{array}{cc} 
\text{Cov}_{\text{From}} & \text{Cov}_{\text{FromCrossTo}} \\ \text{Cov}_{\text{FromCrossTo}}^T & \text{Cov}_{\text{To}}
\end{array} \right]
\end{equation}



The Cholesky decomposition \cite{golub1996matrix} of $\text{Cov}_x$ is computed:
\begin{equation}
\rtc{\text{Cov}_x} = \text{chol}\left( \text{Cov}_x \right)
\end{equation}
where
\begin{equation} \label{eq:chol_decomp}
\text{Cov}_x = \rtc{\text{Cov}_x} \rtc{\text{Cov}_x}^T
\end{equation}

The Cholesky decomposition matrix is notionally considered the ``square root'' of the covariance matrix.  However, the covariance when working with derived points can often be positive semi-definite instead of positive definite: some of the eigenvalues of the matrix are zero as a result of having one or more position components fixed.  The Cholesky decomposition process requires a positive definite matrix input in order to provide a unique solution. Fortunately, the Cholesky decomposition process is stable as the eigenvalues of a matrix approach zero (i.e. changing the eigenvalue from one very small value to another very small value does significantly not change the decomposition), and thus small values can be substituted for any zero-value eigenvalues to obtain a Cholesky decomposition that is correct to any desired precision \cite{wiki:Cholesky,fang2008modified}.

To determine if there are any zero-value eigenvalues before performing the Cholesky decomposition, the eigenvalues and associated eigenvectors for the $\text{Cov}_x$ matrix are computed via the EigenSolver tool within the Eigen library \cite{EigenSolver}:
\begin{equation} \label{eq:eigDecomp}
\text{Cov}_x = U \Lambda U^T
\end{equation}
where $\Lambda$ is a diagonal matrix with eigenvalues on the diagonal, and $U$ is an orthogonal matrix consisting of the eigenvectors associated with the eigenvalues in $\Lambda$.  If any eigenvalues are zero or slightly negative (likely resulting from numerical precision limitations), these values are changed to a very small positive value $\epsilon$.  After this change, the covariance matrix can be reassembled via equation \ref{eq:eigDecomp}.  Then the Cholesky decomposition can be executed to obtain the ``square root'' covariance matrix described in equation \ref{eq:chol_decomp}.  

The covariance square root matrix is then employed to deterministically generate a set of vectors called ``sigma point vectors'':
\begin{equation}
\begin{split}
\mathcal{X} = \left[ \right. & x + \sqrt{n} \rtc{\text{Cov}_x}(:,1),x + \sqrt{n} \rtc{\text{Cov}_x}(:,2), ... ,  \\
                   & \left.x - \sqrt{n} \rtc{\text{Cov}_x}(:,1),x - \sqrt{n} \rtc{\text{Cov}_x}(:,2), ... \right]
\end{split}
\end{equation}
where $n$ is the number of dimensions of the initial vector (six in this two-station example), and the notation $(:,i)$ indicates the $i^{\text{th}}$ column of the $\rtc{\text{Cov}_x}$ matrix. Each column of $\mathcal{X}$ represents a sigma point vector, also described as $\mathcal{X}_i$ for column $i$. 

These sigma point vectors can be used to map the formal covariance of the station positions to the uncertainty in the ellipsoidal distance. Each sigma point vector $\mathcal{X}_i$ is mapped through the nonlinear function $g(x,a,f)$ which includes the conversion from Cartesian ECEF XYZ positions of both stations to geodetic coordinates and applying Vincenty's Method (section \ref{sect:Vincenty}) to compute the ellipsoidal distance:
\begin{equation}
\mathcal{Y}_i = g(\mathcal{X}_i,a,f)
\end{equation}

The mean of the mapped sigma points, which does not necessarily equal the computed ellipsoidal distance using the nominal station positions, is
\begin{equation}
\hat{y} = \frac{1}{2n} \sum_{i=1}^{2n} \mathcal{Y}_i
\end{equation}
The mean and mapped sigma point vectors are used to compute the mapped covariance matrix:
\begin{equation}
\text{Cov}_y = \frac{1}{2n} \sum_{i=1}^{2n} \left(\mathcal{Y}_i - \hat{y}\right) \left(\mathcal{Y}_i - \hat{y}\right)^T
\end{equation}
Note that the dimension of $\hat{y}$ and $\text{Cov}_y$ will be $1$ and $1 \times 1$ for the single geodesic distance output, and thus
\begin{equation}
\sigma_{\text{ellipsoidal}} = \sqrt{\text{Cov}_y}
\end{equation}

% Removing to avoid confusion
% \textbf{Note}: The scaling of the ellipsoid is not included in the sigma point transform, under the assumption it is desired to have a constant ellipsoid based on the nominal ``From'' station position, and not a new ellipsoid for each sigma point.  Additionally, performing this adjustment for each sigma point is not likely to make a significant difference on the mapped uncertainty, while introducing more computation required for each sigma point transformation.


\subsection{Horizontal Distance}

The horizontal distance between two stations is defined as the magnitude of the 2D vector between the stations that results from projecting the relative 3D vector into the horizontal plane.  Given the two station positions in ECEF cartesian frame coordinates, the relative station vector is
\begin{equation} \label{eq:RelVect}
\boldsymbol{r}_{\text{rel}}^{\text{XYZ}} = \boldsymbol{r}_{\text{To}} - \boldsymbol{r}_{\text{From}}
\end{equation}
This $\boldsymbol{r}_{\text{rel}}^{\text{XYZ}}$ vector is projected into the horizontal plane, defined as perpendicular to the WGS 84 ellipsoid at the ``from'' station coordinates:
\begin{equation}
\boldsymbol{r}_{\text{rel}}^{\text{ENU}} = R_{\text{XYZ}}^{\text{ENU}} \boldsymbol{r}_{\text{rel}}^{\text{XYZ}}
\end{equation}
where $R_{\text{XYZ}}^{\text{ENU}}$ is the rotation matrix from the ECEF frame to the ENU frame (which is also fixed to the surface of the Earth). The $R_{\text{XYZ}}^{\text{ENU}}$ rotation matrix is defined as:
\begin{equation}
R_{\text{XYZ}}^{\text{ENU}} = \left[ \hat{E} \: \hat{N} \; \hat{U} \right]^T
\end{equation}
where the unit vectors $\hat{E}$, $\hat{N}$, and $\hat{U}$ are the unit vectors of the east, north, and up directions in the ECEF frame at the ``From'' station. The terms of $R_{\text{XYZ}}^{\text{ENU}}$ are defined in equation \ref{eq:ECEF2ENU}. 

To obtain the horizontal distance, the RSS magnitude of the east and north components is computed:
\begin{equation} \label{eq:HorizontalDist}
D_{\text{horizontal}} = \sqrt{\boldsymbol{r}_{\text{rel}}^{\text{ENU}}(1)^2 + \boldsymbol{r}_{\text{rel}}^{\text{ENU}}(2)^2}
\end{equation}

To determine the uncertainty in the horizontal distance, the same unscaled sigma point transform is employed as described above for the ellipsoidal distance. The mapping function is the horizontal distance calculation provided in equations \ref{eq:RelVect} through \ref{eq:HorizontalDist} instead of the ellipsoidal distance, but otherwise the approach is identical.
